\chapter{Личность: не таблица весов}
\label{ch:personality}

\section{Главная ловушка}

Соблазнительное, но неверное решение — «слой личности» как хардкод: таблица весов и
жёстких триггеров поверх графа («умножь ребро злорадства на $2{,}0$, обрежь сочувствие
до $0{,}05$, на узел „грусть“ повесь триггер инверсии $\to$ сарказм»).

\begin{pitfall}[Почему таблица весов — тупик]
\emph{Во-первых}, это стек правил, противоречащий установке ядра («черты — следствие
мотиваций, а не правила поверх»). \emph{Во-вторых}, это карикатура: таблица весов
даёт \emph{паттерновое}, но \emph{не неожиданное} и одноголосое поведение — всегда
сарказм, потому что так прошито; а живая личность — неожиданные, но паттерновые
реакции. \emph{В-третьих}, это ровно ошибка жёсткой скалярной декомпозиции: реакция
на «разбил кружку» зависит от настроения, отношений и всплывшей памяти; фиксированное
ребро «грусть$\to$сарказм» это выкидывает.
\end{pitfall}

\begin{workedbox}[Формально: почему скалярный роутинг теряет взаимодействие]
Пусть желаемый ответ — функция контекста $y=\Psi(\mathbf{g},\mathbf{l})$, где
$\mathbf{g}$ — глобальная аранжировка (настроение, отношения), $\mathbf{l}$ —
локальный триггер (событие). Таблица весов реализует $y\approx a\cdot\mathbf{l}$ —
ранг-1 скалярный роутинг, не зависящий от $\mathbf{g}$. Но истинное $\Psi$ содержит
\emph{перекрёстные} члены $\mathbf{g}^\top W \mathbf{l}$ (реакция на событие
\emph{зависит} от аранжировки). Ранг-1 приближение их обнуляет — теряется ровно то,
что делает характер контекстным. Это тот же провал, что жёсткая coarse-plus-local
декомпозиция в геометрических графовых моделях.
\end{workedbox}

\section{Личность распределена}

Личность \emph{не} слой весов в графе. Она распределена по трём источникам:
\emph{мотивационная структура} (мастерство, участие, самоуважение) $+$ \emph{текущее
настроение} $+$ \emph{значимость всплывшего из памяти}. И всё это стилизует
\emph{трансформ Rinness} (голос, глава~\ref{ch:voice}).

На стенде это — не статические множители, а \emph{смещение во время обхода}:
настроение и мотивация непрерывно модулируют переходы \eqref{eq:trans}. Тогда
грустная и злая Рин на одно событие всплывают разное — характер контекстно-зависим.
Держи веса на рёбрах для \emph{силы связи и уверенности}; личность — как смещение от
состояния, не как sarcasm-множитель.

\section{Граф хранит ингредиенты, а не реакцию}

Возражение: «если из „разбить кружку“ ведёт только „грусть“, другую реакцию взять
негде». Ответ в двух частях. Во-первых, база причинно-социального смысла (ATOMIC,
глава~\ref{ch:experiment}) даёт не одно ребро, а \emph{пучок}: намерение, реакция,
эффект, реакция окружающих. Настроение смещает ветку пучка. Во-вторых, и главное:
\emph{реакцию не берут ребром — её порождают}. Обход собирает активированную
окрестность-подграф, а реакцию Rinness/LLM \emph{композирует поверх} материала под
настроение и мотивацию. Настроение входит \emph{дважды}: на всплывании (грустная
поднимает «потеря», игривая — «неуклюжесть») и на композиции.

\section{Консолидация реакций: по значимости, а не по приёму}

\subsection{Провизорная запись}

Сгенерированная реакция пишется провизорным ребром с низкой уверенностью (кандидат;
принцип «гипотеза, а не факт»). Ответ «это было лишнее» — отрицательное
\emph{свидетельство}, обновляющее уверенность того же ребра вниз, а не конкурирующая
консолидация. Механически ребро-реакция несёт счётчик успехов/провалов
$\mathrm{Beta}(\alpha,\beta)$:
\begin{equation}
\text{успех}:\ \alpha\!\mathrel{+}\!=\!1;\quad \text{отпор}:\ \beta\!\mathrel{+}\!=\!1;\quad
\conf=\frac{\alpha}{\alpha+\beta},\ \ \text{неопред.}=\mathrm{Var}[\mathrm{Beta}].
\end{equation}
Стабильно принимается — переходит порог и консолидируется в медленный слой;
стабильно ловит отпор — угасает. Выбор реакции по такому апостериору — это
сэмплирование Томпсона: баланс исследования новых реакций и эксплуатации
проверенных.

\subsection{Гейт значимости, а не одобрения}

\begin{pitfall}[Нельзя гейтить консолидацию приёмом]
Консолидировать реакции по тому, как их приняли, значит строить угодника, чья
личность формируется одобрением, — ровно то, что запрещает ядро («самоуважение не
течёт через понравиться»). Гейт характера внешним приёмом философски кривой для Рин.
\end{pitfall}

Правильный ход — развернуть гейт \emph{внутрь}: консолидируй по \emph{внутренней
значимости} (само-оценке при кодировании). Что такое значимость \emph{конкретно}?
Это \emph{теория оценивания} (appraisal: OCC; Шерер). Событие оценивается по
измерениям — surprise (ошибка предсказания), goal-congruence (скалярное произведение
с вектором мотивации), valence, intensity, agency, — и значимость
\begin{equation}
\text{salience}\;\propto\;\text{surprise}\times\text{goal\_relevance}\;(+\,\text{intensity}).
\label{eq:salience}
\end{equation}
Два замечания. \emph{Дешевизна}: простейшая конкретизация — $|\Delta\mood|$,
величина изменения настроения, а она уже посчитана; mood-модуль работает детектором
значимости почти без лишних вычислений. \emph{Глубина}: самый глубокий одиночный
сигнал — ошибка предсказания, и она же питает \texttt{need\_to\_discover}. Значимость
и обучение — один сигнал.

Реакцию собеседника Рин помнит и оценивает — но как \emph{просто событие}: с низкой
конгруэнтностью её ценностям «ему не зашло» обновляет модель отношения и чуть-чуть
настроение, но не доминирует в консолидации характера. Рамка оценивания сама даёт
«считаться, а можно и нет, частично».

\subsection{Условие по отношениям}

Одна шутка заходит собеседнику $A$ и падает у $B$: ребро-реакция квалифицируется
парой (ситуация $\times$ отношение), консолидируется по-персонно. Один негативный
акт даёт три эффекта: понижение ребра, апдейт настроения (осеклась),
мета/relational-пометку.

\begin{workedbox}[«Разбил кружку», полный проход]
(1)~Событие активирует окрестность: «кружка», «грусть», «неуклюжесть», плюс личная
память «$B$ хотел новую» ($-$ contradicts «грусть»). (2)~Настроение (игривое)
смещает всплывание к «неуклюжесть». (3)~Rinness композирует подкол — не берёт ребром,
собирает из ингредиентов. (4)~Ребро-реакция пишется провизорно,
$\mathrm{Beta}(1,1)$. (5)~$B$ смеётся — $\alpha\!+\!1$; или «это лишнее» —
$\beta\!+\!1$. (6)~Гейт консолидации — по $|\Delta\mood|$ и surprise, не по одобрению
$B$; но реакция $B$ обновляет модель «$B\times$такой юмор». После многих проб
подкол-на-неуклюжесть у $B$ либо консолидируется, либо угасает — комедийный тайминг
выучен как Бета-апостериор.
\end{workedbox}

\section{Дистиллят как система координат}

Возникает соблазн вовсе убрать «дистиллят» (усреднённый здравый смысл готовых баз).
Не надо. Дистиллят — \emph{общая модель мира}, относительно которой отклонение
\emph{читаемо} как характер: личность «неожиданна» лишь на фоне общих ожиданий; без
общей почвы причуды читаются как шум. Правильно: дистиллят фоном (система координат),
личность — процесс порождения отклонений на нём. Без базового «разбитая вещь обычно
плохо» не понять, почему здесь «поздравляю» — уместный сарказм.
